% !TeX program = xelatex
% ============================================================
%  Software Requirements Specification (SRS)
%  ISO/IEC/IEEE 29148-2018 §9.6
%
%  English international-standard ANALOG of the ESPD Statement of
%  Work (GOST 19.201-78). Per §2.3.4.2 of the HSE SE practice
%  programme, a student writing in English may submit this SRS
%  instead of the translated GOST 19 form.
%
%  Self-contained, compile-ready xelatex document. English only:
%  there is no Russian branch and no project.lang conditionals.
% ============================================================
\documentclass[12pt,a4paper]{extarticle}
\newcommand{\hseDefaultMono}{Consolas}
\input{../common/fonts}
\usepackage[english]{babel}
\usepackage[a4paper,left=25mm,right=15mm,top=20mm,bottom=20mm]{geometry}
\usepackage{setspace}\onehalfspacing
\usepackage{indentfirst}\setlength{\parindent}{1.25cm}\setlength{\parskip}{6pt}
\usepackage{microtype}\usepackage{csquotes}\usepackage{amsmath}\usepackage{graphicx}
\usepackage{xcolor}\usepackage{float}\usepackage{array}\usepackage{booktabs}
\usepackage{longtable}\usepackage{tabularx}\usepackage{adjustbox}\usepackage{xurl}\urlstyle{same}
\usepackage{hyphenat}\usepackage{caption}\usepackage{enumitem}\setlist{nosep}
\usepackage{titlesec}\usepackage{hyperref}
\hypersetup{unicode=true,colorlinks=true,linkcolor=black,urlcolor=blue}
\input{../common/meta}    % provides \hseFill, \projectname, \hseProjectNameEn, \hseAuthorName, \hseGroup, \hseYear ...
\input{../common/commands}

% ── Local table column types ────────────────────────────────
\newcolumntype{L}[1]{>{\raggedright\arraybackslash}p{#1}}

\begin{document}\sloppy

% ============================================================
% TITLE BLOCK
% ============================================================
\begin{center}
{\large\bfseries Software Requirements Specification\par}
\vspace{2pt}
{\normalsize according to ISO/IEC/IEEE 29148-2018\par}
\vspace{16pt}
{\large\bfseries \ifhseHasNameEn\hseProjectNameEn\else\projectname\fi\par}
\vspace{18pt}
{\normalsize
Author: \hseAuthorName\par
Group: \hseGroup\par
Year: \hseYear\par}
\end{center}

\vspace{6pt}

\begin{abstract}
\noindent This document specifies the software requirements for the system
\enquote{\ifhseHasNameEn\hseProjectNameEn\else\projectname\fi}. It is
structured in accordance with ISO/IEC/IEEE 29148-2018 (Systems and software
engineering~--- Life cycle processes~--- Requirements engineering), clause~9.6,
and is submitted as the English international-standard equivalent of the
Statement of Work. Each requirement is identified so that it can be traced to
design, implementation and verification artefacts.
\end{abstract}

\hypersetup{linkcolor=black}
\tableofcontents
\hypersetup{linkcolor=blue}
\clearpage

% ============================================================
% 1. INTRODUCTION
% ============================================================
\section{Introduction}

\subsection{Purpose}
\label{sec:purpose}

This Software Requirements Specification (SRS) defines the functional and
non-functional requirements for \enquote{\ifhseHasNameEn\hseProjectNameEn\else\projectname\fi}.
It is intended for \hseFill{intended readers, e.g. developers, reviewers,
supervisor} and serves as the agreed basis for design, implementation,
verification and acceptance of the software product.

\subsection{Scope}
\label{sec:scope}

The software product to be produced is named
\enquote{\ifhseHasNameEn\hseProjectNameEn\else\projectname\fi}.
The system \hseFill{what the system will do}. It will not
\hseFill{explicit out-of-scope items}. The expected benefits and objectives
are \hseFill{main goals and benefits}.

\subsection{Product overview}

\subsubsection{Product perspective}

The system is \hseFill{self-contained product / component of a larger system}.
It interacts with \hseFill{external systems, services or actors} through the
interfaces defined in section~\ref{sec:external_interfaces}. A high-level
context of the system is shown in \hseFill{reference to a context diagram}.

\subsubsection{Product functions}

At a high level, the system provides the following functions:
\begin{itemize}[leftmargin=1.25cm]
    \item \hseFill{principal function 1};
    \item \hseFill{principal function 2};
    \item \hseFill{principal function 3}.
\end{itemize}
The detailed functional requirements are specified in
section~\ref{sec:functional_requirements}.

\subsubsection{User characteristics}

The intended classes of users, together with their expected expertise and
responsibilities, are:
\begin{itemize}[leftmargin=1.25cm]
    \item \textbf{\hseFill{user class 1, e.g. end user}}~--- \hseFill{expertise and tasks};
    \item \textbf{\hseFill{user class 2, e.g. administrator}}~--- \hseFill{expertise and tasks}.
\end{itemize}

\subsubsection{Assumptions and dependencies}

The following assumptions and dependencies affect the requirements in this
document:
\begin{itemize}[leftmargin=1.25cm]
    \item \hseFill{assumption, e.g. availability of an external service};
    \item \hseFill{dependency, e.g. specific runtime or third-party library}.
\end{itemize}

\subsection{Definitions and acronyms}
\label{sec:definitions}

\begin{longtable}{L{3.5cm} L{\dimexpr\textwidth-3.5cm-2\tabcolsep\relax}}
\toprule
\textbf{Term / acronym} & \textbf{Definition} \\
\midrule
\endfirsthead
\toprule
\textbf{Term / acronym} & \textbf{Definition} \\
\midrule
\endhead
\bottomrule
\endfoot
API & Application Programming Interface. \\
SRS & Software Requirements Specification. \\
\hseFill{term} & \hseFill{definition}. \\
\hseFill{acronym} & \hseFill{expansion and definition}. \\
\end{longtable}

\clearpage

% ============================================================
% 2. REFERENCES
% ============================================================
\section{References}
\label{sec:references}

\begin{enumerate}[label={[}\arabic*{]}, leftmargin=1.25cm]
    \item ISO/IEC/IEEE 29148-2018, \enquote{Systems and software engineering~--- Life cycle processes~--- Requirements engineering.} International Organization for Standardization, 2018.
    \item \hseFill{project charter or assignment reference}.
    \item \hseFill{applicable external standard or specification}.
    \item \hseFill{related design or interface document}.
\end{enumerate}

\clearpage

% ============================================================
% 3. SPECIFIC REQUIREMENTS
% ============================================================
\section{Specific requirements}

\subsection{External interface requirements}
\label{sec:external_interfaces}

\subsubsection{User interfaces}

\hseFill{describe the user interface: screens, layout conventions, input/output
elements, accessibility}. The user interface shall \hseFill{key UI requirement,
e.g. present validation errors inline}.

\subsubsection{Hardware interfaces}

\hseFill{describe hardware interfaces, or state that the system has no direct
hardware interfaces}. The system shall run on \hseFill{target hardware or virtual
environment}.

\subsubsection{Software interfaces}

The system interacts with the following software components:
\begin{itemize}[leftmargin=1.25cm]
    \item \textbf{\hseFill{component, e.g. database}}~--- \hseFill{purpose, version, protocol};
    \item \textbf{\hseFill{external service or API}}~--- \hseFill{purpose, version, data exchanged}.
\end{itemize}

\subsubsection{Communications interfaces}

\hseFill{describe network protocols, ports, data formats and security of
communication, e.g. HTTPS + JSON}. Data exchanged over the network shall be
encoded as \hseFill{data format, e.g. JSON (UTF-8)}.

\subsection{Functional requirements}
\label{sec:functional_requirements}

Each functional requirement is identified by a unique label
(\texttt{SRS-F-01}, \texttt{SRS-F-02}, \ldots) so that it can be traced to
design, implementation and the verification cases in
section~\ref{sec:verification}.

\begin{longtable}{L{2.2cm} L{\dimexpr\textwidth-6.7cm-4\tabcolsep\relax} L{1.8cm} L{2.7cm}}
\caption{Functional requirements}\label{tab:functional_reqs}\\
\toprule
\textbf{ID} & \textbf{Requirement} & \textbf{Priority} & \textbf{Source} \\
\midrule
\endfirsthead
\multicolumn{4}{l}{\small\itshape Table~\ref{tab:functional_reqs} (continued)}\\
\toprule
\textbf{ID} & \textbf{Requirement} & \textbf{Priority} & \textbf{Source} \\
\midrule
\endhead
\midrule
\multicolumn{4}{r}{\small\itshape continued on the next page}\\
\endfoot
\bottomrule
\endlastfoot
SRS-F-01 & The system shall \hseFill{observable behaviour, e.g. allow a user to register}. & \hseFill{High} & \hseFill{use case / stakeholder} \\
SRS-F-02 & The system shall \hseFill{observable behaviour}. & \hseFill{Medium} & \hseFill{source} \\
SRS-F-03 & The system shall \hseFill{observable behaviour}. & \hseFill{Low} & \hseFill{source} \\
\end{longtable}

\subsection{Usability requirements}
\label{sec:usability}

\begin{itemize}[leftmargin=1.25cm]
    \item \textbf{SRS-U-01}~--- A \hseFill{user class} shall be able to complete \hseFill{key task} within \hseFill{measurable target, e.g. 3 steps / 2 minutes}.
    \item \textbf{SRS-U-02}~--- Error messages shall \hseFill{usability criterion, e.g. state the cause and a corrective action}.
    \item \textbf{SRS-U-03}~--- The interface shall conform to \hseFill{style guide or accessibility standard}.
\end{itemize}

\subsection{Performance requirements}
\label{sec:performance}

\begin{itemize}[leftmargin=1.25cm]
    \item \textbf{SRS-P-01}~--- The system shall process \hseFill{operation} within \hseFill{time, e.g. 500 ms} for the 95th percentile of requests.
    \item \textbf{SRS-P-02}~--- The system shall support at least \hseFill{number} concurrent users under normal load.
    \item \textbf{SRS-P-03}~--- The system shall sustain a throughput of \hseFill{requests per second and justification}.
\end{itemize}

\subsection{Logical database requirements}
\label{sec:database}

The persistent data managed by the system is described below.
\begin{itemize}[leftmargin=1.25cm]
    \item \textbf{Entities}: \hseFill{principal entities and their key attributes};
    \item \textbf{Relationships}: \hseFill{relationships and cardinalities};
    \item \textbf{Integrity}: \hseFill{integrity and consistency constraints};
    \item \textbf{Retention}: \hseFill{data retention and archival requirements}.
\end{itemize}

\subsection{Design constraints}
\label{sec:constraints}

\begin{itemize}[leftmargin=1.25cm]
    \item Programming languages and versions: \hseFill{languages and versions};
    \item Frameworks and libraries: \hseFill{mandated frameworks and libraries};
    \item Standards to be observed: \hseFill{coding / regulatory standards};
    \item Deployment environment: \hseFill{runtime and deployment constraints}.
\end{itemize}

\subsection{Software system attributes}
\label{sec:attributes}

\subsubsection{Reliability}
The system shall \hseFill{reliability requirement, e.g. recover to a consistent
state after a component failure without data loss}.

\subsubsection{Availability}
The system shall provide an availability of no less than \hseFill{target, e.g.
99.5\%} over \hseFill{measurement period}, excluding scheduled maintenance.

\subsubsection{Security}
\begin{itemize}[leftmargin=1.25cm]
    \item Authentication: \hseFill{authentication mechanism};
    \item Authorization: \hseFill{access-control model};
    \item Data protection: \hseFill{encryption in transit and at rest};
    \item Auditing: \hseFill{security-relevant events to be logged}.
\end{itemize}

\subsubsection{Maintainability}
The system shall \hseFill{maintainability requirement, e.g. be organised into
independently deployable modules with automated tests}.

\subsubsection{Portability}
The system shall \hseFill{portability requirement, e.g. run on any host that
provides the specified container runtime}.

\clearpage

% ============================================================
% 4. VERIFICATION
% ============================================================
\section{Verification}
\label{sec:verification}

This section states how each class of requirement is verified. The accepted
methods are inspection (I), analysis (A), demonstration (D) and test (T). Each
verification case is identified (\texttt{TC-01}, \texttt{TC-02}, \ldots) and
references the requirement it verifies.

\begin{longtable}{L{1.8cm} L{2.6cm} L{\dimexpr\textwidth-9.4cm-4\tabcolsep\relax} L{2.4cm} L{2.6cm}}
\caption{Verification cases}\label{tab:verification}\\
\toprule
\textbf{TC ID} & \textbf{Verifies} & \textbf{Method / criterion} & \textbf{Method} & \textbf{Expected result} \\
\midrule
\endfirsthead
\multicolumn{5}{l}{\small\itshape Table~\ref{tab:verification} (continued)}\\
\toprule
\textbf{TC ID} & \textbf{Verifies} & \textbf{Method / criterion} & \textbf{Method} & \textbf{Expected result} \\
\midrule
\endhead
\midrule
\multicolumn{5}{r}{\small\itshape continued on the next page}\\
\endfoot
\bottomrule
\endlastfoot
TC-01 & SRS-F-01 & \hseFill{steps or acceptance criterion} & \hseFill{T} & \hseFill{expected outcome} \\
TC-02 & SRS-P-01 & \hseFill{load profile and measurement} & \hseFill{T} & \hseFill{p95 within target} \\
TC-03 & SRS-U-01 & \hseFill{usability observation} & \hseFill{D} & \hseFill{task completed within target} \\
TC-04 & SRS-A-security & \hseFill{security check} & \hseFill{A} & \hseFill{no unauthorized access} \\
\end{longtable}

\noindent Verification of requirement classes:
\begin{itemize}[leftmargin=1.25cm]
    \item \textbf{Functional} (section~\ref{sec:functional_requirements})~--- verified by test against the cases in Table~\ref{tab:verification}.
    \item \textbf{Usability} (section~\ref{sec:usability})~--- verified by demonstration with representative users.
    \item \textbf{Performance} (section~\ref{sec:performance})~--- verified by test under the declared load profile.
    \item \textbf{Software system attributes} (section~\ref{sec:attributes})~--- verified by \hseFill{analysis / inspection / test as applicable}.
\end{itemize}

\clearpage

% ============================================================
% 5. SUPPORTING INFORMATION / APPENDICES
% ============================================================
\section{Supporting information}
\label{sec:supporting}

\subsection*{Appendix A. Traceability}
\addcontentsline{toc}{subsection}{Appendix A. Traceability}

Each requirement identifier (\texttt{SRS-F-*}, \texttt{SRS-U-*},
\texttt{SRS-P-*}) is traced to its verification case (\texttt{TC-*}) in
Table~\ref{tab:verification}. \hseFill{reference to the full traceability matrix,
if maintained separately}.

\subsection*{Appendix B. Open issues}
\addcontentsline{toc}{subsection}{Appendix B. Open issues}

\begin{itemize}[leftmargin=1.25cm]
    \item \hseFill{open issue or decision to be confirmed};
    \item \hseFill{item requiring further analysis}.
\end{itemize}

\end{document}
